\section{Additional steps in cote 149, scans 41--104}
\label{sec:ZCF}

\begin{result}{ZCF1}{Known}{The open birational Hom statement, and the withdrawn unrestricted version}
\textup{[\S2, manuscript pp.\ 20--22; scans 41--46.]}
Let $k$ be a number field, and let $L,M$ be finitely generated regular
extensions of $k$. With compatible separable closures, restriction gives
a bijection
\[
 \operatorname{Hom}_{k\text{-fields}}(M,L)
 \longrightarrow
 \operatorname{Hom}^{\mathrm{open}}_{G_k}(G_L,G_M)
             /\operatorname{Inn}(G_{M\bar k}).                 \tag{ZCF.1}
\]
All homomorphisms on the right are continuous. Thus a) and b) of
scan 44 are equivalent after the common constant field has been fixed
and the geometric kernels are the kernels for that same field.
No supplementary compatibility with a filtration by models is required
for this open-Hom statement. The version with arbitrary homomorphisms
in scans 41--42 is false, as the manuscript itself observes in scan 43.
\end{result}
\begin{proof}
The precise input is Mochizuki, Theorem 17.1
\cite{CF_Mochizuki1999}: the ground field is sub-$p$-adic, the two
function fields have arbitrary finite transcendence degree, and the
ground field is algebraically closed in each, as specified immediately
before the theorem. A number field is sub-$p$-adic via any place over
$p$, so these hypotheses hold. The theorem is exactly (ZCF.1), with
the displayed geometric inner-conjugacy ambiguity.

For clarity, the easy direction has its own field-theoretic proof.
For $M\hookrightarrow L$, the algebraic closure $M_0$ of $M$ in $L$
is finite over $M$, and $L/M_0$ is regular. Linear disjointness from
$M_0^{\mathrm{sep}}$ makes $G_L\to G_{M_0}$ surjective. Its composite
into $G_M$ therefore has finite-index image $G_{M_0}$.
If $\phi:E'\to E$ is over the identity of $G_k$, with geometric
kernels $P',P$, then
\[
 \phi(E')\cap P=\phi(P'),\qquad
 [E:\phi(E')]=[P:\phi(P')].
\]
The equality of intersections follows by applying the quotient map;
every coset of $\phi(E')$ has a representative in $P$, since
$\phi(E')\to G_k$ is surjective. This proves the equality of indices
and the manuscript's arithmetic/geometric openness comparison, with
the directions of the arrows corrected.

Take instead $L=k$ and $M=k(t)$. The coefficient action on compatible
Puiseux roots at $t=0$ gives a continuous section $G_k\to G_{k(t)}$.
It is a morphism over $G_k$, whereas no $k$-field embedding
$k(t)\to k$ exists. The corresponding map of geometric groups is
$1\to G_{\bar k(t)}$, so this also refutes the proposed unrestricted
geometric-Hom version. The coefficient construction is written out in
F9.3. Its image has infinite index, since its intersection with the
infinite geometric kernel is trivial.

Finally $G_M=\varprojlim_W\pi_1(W)$, where $W$ runs through nonempty
opens of a fixed normal finite-type $k$-model of $M$. A finite separable
extension of $M$ has finite normalization on that model and is finite
\'etale after deleting its branch locus; the same statement for a
compositum and for morphisms of finite extensions proves the equality
of Galois categories and then the displayed limit. A8 supplies cofinal
elementary models. This proves the limit assertion in scan 46 without
using a conjectural preservation of centralizers by inverse limits.
\end{proof}

\begin{result}{ZCF2}{Proved}{Recovering an extension map from an equivariant geometric map}
\textup{[\S2, manuscript p.\ 21; scans 43--44.]}
Let $1\to P_i\to E_i\to G\to1$ be profinite extensions and let
$\phi:P_1\to P_2$ be continuous. Suppose its conjugacy class is
$G$-equivariant and $C_{P_2}(\phi(P_1))=1$. Then $\phi$ extends
uniquely to a continuous homomorphism $F:E_1\to E_2$ over $G$.
Here equivariance means that, for every $e_1\in E_1$ and any $e_2\in E_2$
with the same image in $G$, the maps
$\phi\circ\operatorname{int}(e_1)$ and
$\operatorname{int}(e_2)\circ\phi$ differ by an inner automorphism
from $P_2$. The stronger hypothesis involving every open subgroup in
scan 43 implies the centralizer hypothesis just stated.
\end{result}
\begin{proof}
For $e_1$ choose $e_2$ as above. Equivariance supplies $b\in P_2$
such that $be_2$ satisfies
\[
 \operatorname{int}(be_2)\circ\phi
       =\phi\circ\operatorname{int}(e_1).
\]
There is only one such element above the given element of $G$: the
ratio of two would centralize $\phi(P_1)$. Define $F(e_1)=be_2$.
Uniqueness proves both the multiplication law and $F|_{P_1}=\phi$.
Its graph is the closed subset of $E_1\times E_2$ cut out by equality
in $G$ and all the equations
$e_2\phi(p)e_2^{-1}=\phi(e_1pe_1^{-1})$, $p\in P_1$.
Projection of this compact graph to $E_1$ is a continuous bijection
to a Hausdorff space, hence a homeomorphism. Thus $F$ is continuous.
Changing $\phi$ by a $P_2$-inner conjugation changes $F$ by that same
conjugation, specifying the outer ambiguity.

The tentative arithmetic centralizer condition c) of scan 44 does
not replace openness. For the coefficient section at $t=0$ above,
the attached return's B7 computes the centralizer of the cyclotomic
kernel as the local inertia $T$. For an open subgroup of the full
$G_k$, its centralizer is contained in $T$, and invariance there is
$\chi(g)a=a$. Each $\ell$-adic cyclotomic character has infinite
image over a number field, so $a_\ell=0$ for every $\ell$.
Projection to $G_k$ adds no centralizer, by F8a. Thus all these
arithmetic centralizers are trivial although the image is not open.
This also refutes the withdrawn assertion $P^{G_{k'}}\ne1$.
The corrected cyclotomic assertion $P^{G_{k'}^\circ}\ne1$ for an
\emph{arbitrary} birational section is retained as Open in ZCF6.
\end{proof}

\begin{result}{ZCF3}{Refuted}{Geometric homology alone does not determine a morphism}
\textup{[\S2, manuscript p.\ 23; scan 47.]}
Two morphisms into an elementary hyperbolic variety can induce the
same map on every geometric $H_1(-,\mathbb Z_\ell)$ and different
arithmetic fundamental-group homomorphisms. Consequently the
unqualified sentence in scan 47 is false, for either direction of
the doubtful source/target notation.
\end{result}
\begin{proof}
Set $Y=\mathbb P^1_{\mathbb Q}\setminus\{0,1,\infty\}$ and take
the two morphisms $\operatorname{Spec}\mathbb Q\to Y$ with values
$2$ and $3$. Both geometric homology maps are zero. The unit $x$ on
$Y$ induces $\pi_1(Y)\to\pi_1(\mathbb G_m)$ over $G_{\mathbb Q}$.
On the two sections it gives the Kummer classes of $2$ and $3$.
Their difference is the class of $3/2$; its $2$-adic valuation is
$-1$, so it is nonzero in
$\mathbb Q^*/\mathbb Q^{*\ell^n}$ for every $\ell,n>0$.
Thus their arithmetic sections differ even in each geometric
pro-$\ell$ abelian quotient. To have a positive-dimensional source,
compose these constants with $Y\to\operatorname{Spec}\mathbb Q$;
restriction at a rational point gives the same distinction.
The proved replacements are the arithmetic Kummer injectivity B2
and the dominant-map rigidity B5. The displayed quotient maps explain
exactly which information the geometric-homology sentence loses.
\end{proof}

\begin{result}{ZCF4}{Refuted}{An explicit nondominant map inducing an isomorphism of full fundamental groups}
\textup{[\S2, manuscript pp.\ 24--25; scans 49--50.]}
Let $C/\mathbb Q$ be the smooth projective completion of
$y^2=x^5-x$, and let $Y=C^3$. On $(\mathbb P^1)^3$ take the
trihomogeneous polynomial whose affine expression is
\[
 F=1+7x_1+3x_2+2x_3+13x_1x_2+11x_1x_3
                     +5x_2x_3+19x_1x_2x_3.
\]
Let $H\subset Y$ be its inverse-image zero divisor under the product
of the three hyperelliptic maps. Then $H$ is a smooth geometrically
connected surface and the closed immersion $H\hookrightarrow Y$
induces an isomorphism $\pi_1(H)\simeq\pi_1(Y)$ over $G_{\mathbb Q}$.
It is not dominant. Thus the open-Hom assertion of scan 50 is false
for the stated higher-dimensional elementary targets, already for
constant inverse systems.
\end{result}
\begin{proof}
The six branch values of $C\to\mathbb P^1$ are
$B=\{0,1,-1,i,-i,\infty\}$. We verify smoothness over
$\overline{\mathbb Q}$. The eight affine charts of $F=0$ in
$(\mathbb P^1)^3$ have singularity ideals generated by $F$ and its
three first derivatives. Exact rational Gr\"obner calculation gives
the unit ideal in every chart. Fixing any one coordinate to any
element of $B$ leaves a nonsingular bilinear form in the other two
coordinates: all eighteen $2\times2$ determinants are nonzero.
Fixing all three to branch values gives 216 nonzero values.
These exhaustive checks and their exact output are in
\texttt{code/CF\_hyperplane.py} and
\texttt{code/CF\_hyperplane\_output.json}.

Here is why these finite checks prove the geometric assertion.
At a point of $Y$, a local hyperelliptic coordinate is unramified
unless its $x$-coordinate is in $B$; at a branch point it has the
form $x-b=u^2$ up to a unit (with $1/x$ at infinity).
With no ramified coordinate, smoothness follows from the first
check. With one, the nonsingular bilinear restriction makes at
least one of the two unramified derivatives nonzero. With two,
the remaining linear form is nonzero, since a nonsingular
bilinear form cannot vanish identically after fixing one of its
variables to a projective point; its zero has nonzero derivative.
With three, the third check says that the point is not on $H$.
This proves smoothness everywhere.

The pullback of $\mathcal O(1,1,1)$ is ample, because the product
hyperelliptic map is finite. Its section $F$ is nonzero and defines
an effective Cartier divisor in the integral threefold $Y$.
SGA~2 XII, Corollary 3.5 \cite{CF_SGA2} applies: the closed local
rings of $Y$ have depth three, and all local rings off the divisor
are regular, hence pure. It gives geometric connectedness and
$\pi_1(H_{\overline{\mathbb Q}})\simeq
\pi_1(Y_{\overline{\mathbb Q}})$. The two arithmetic exact
sequences have this isomorphism on kernels and the identity on
$G_{\mathbb Q}$, hence give the asserted arithmetic isomorphism.
Dimensions give $\dim H=2<3=\dim Y$.

The failure mechanism is also explicitly stated in the Remark
following Mochizuki's Theorem 17.1. A precise proved replacement
for the one-dimensional-target version is his Theorem 16.5 and
its following Remark: a smooth pro-variety and a hyperbolic
pro-curve over a sub-$p$-adic field, with birational transition
maps, have dominant morphisms in bijection with continuous open
maps of their full arithmetic fundamental groups modulo geometric
inner automorphisms. Its full profinite comparison is proved there
by pulling back every finite cover; it does not claim the false
higher-dimensional-target version.
\end{proof}
\begin{result}{ZCF5}{Proved}{The centralizer precisions and the necessary half of the fundamental conjecture}
\textup{[\S2--3, manuscript pp.\ 23--31; scans 48--49, 53--63.]}
For a hyperbolic curve $U/k$, with $k$ finitely generated over
$\mathbb Q$, and a continuous section $s:G_k\to E_U$, put
$G_k^\circ=\ker\chi$. The checked attached return A6--A7 and AE1 \cite{PriorCyclotomic} gives
\[
 C_\Pi(s(G_k^\circ))=
 \begin{cases}I_c,&s\text{ is cuspidal at its unique lifted cusp }c,\\
               1,&s\text{ is noncuspidal.}\end{cases}          \tag{ZCF.2}
\]
For every open $\Omega\subset G_k^\circ$ the same formula holds
with $\Omega$ in place of $G_k^\circ$. In particular genuine finite
point sections have zero cyclotomic centralizer. The same zero
centralizer conclusion holds for point sections of hyperbolic
polycurves. It proves a)$\Rightarrow$b) in scan 48 for actual
morphisms into these varieties.
\end{result}
\begin{proof}
We use the proved return, not the unsupported sketch R2. Its A6
uses Stix Theorems 14 and 20 \cite{CF_Stix2008}: cyclotomic
normalization detects peripheral subgroups, and the cuspidal packets
are disjoint from each other and from finite-point sections. Its
arithmetic hypothesis over $k(\mu_\infty)$ is supplied by the
proved absence of divisible abelian-variety points in AE1. A7
passes to any open $\Omega$ by writing its fixed field as
$k'(\mu_\infty)$ with $k'/k$ finite, and descends the unique cusp
using normality after a finite Galois enlargement. We inspected
these proofs, including the finite-cover cusp uniqueness step.
Thus (ZCF.2) has precisely the needed hypotheses; it says
noncuspidal, without asserting that such a section comes from a point.

For a polycurve, project a centralizing element through its curve
tower. The base point section has zero centralizer by induction.
The element therefore lies in the fibre kernel; C4 identifies that
kernel with the full geometric group of the fibre through the point.
The last curve case of (ZCF.2) makes it one. The same argument works
after each finite field extension and for each open $\Omega$.

Now let $f:V\to U$ be an actual morphism as in scan 48, with $V$
nonempty and of finite type over a number field, and let
$H\subset E_V^\circ$ be open. Choose a closed point $v$ of $V$.
Its residue field $k(v)$ is finite over the ground number field,
and its point section has image in $E_V$. The inverse image of
$H$ in $G_{k(v)}^\circ$ is open there. The composite point section
in $E_U$ is the section of $f(v)$ after this residue-field base
change, so its geometric centralizer is one by the preceding
paragraph. Any element centralizing $f_*(H)$ centralizes this
smaller image. Its arithmetic projection is also one by F8a.
Therefore $C_{E_U}(f_*(H))=1$, as required.

The lifted-cusp assertion in scan 58 has an explicit group map.
Above a fixed geometric cusp the lifted cusps are $\Pi/I_c$;
the map $aI_c\mapsto aI_ca^{-1}$ is well defined and surjective
onto its conjugates. It is injective because $N_\Pi(I_c)=I_c$
(A4/F10a). Distinct base cusps give disjoint conjugacy classes by
the same peripheral intersection theorem. This is the claimed
recovery of actual lifted cusps, not just their unlabelled orbits.

Finally, for a closed subgroup $\Delta\subset G_k$ containing such
an $\Omega$, the cuspidal case gives the exact formula
\[
 C_\Pi(s(\Delta))=I_c^{\chi(\Delta)}
 =\prod_{\ell:\,\chi_\ell|_\Delta=1}\mathbb Z_\ell(1).
\]
It follows coordinatewise from
$(\chi_\ell(\delta)-1)a_\ell=0$ in the integral domain
$\mathbb Z_\ell$. This is the corrected H1, including the possibility
of nonzero invariants with a nontrivial full cyclotomic character.
\end{proof}

\begin{result}{ZCF6}{Open}{The existence assertions retained after the centralizer computation}
\textup{[\S2--3, manuscript pp.\ 22--34 and 51; scans 45, 48--49,
51, 56--57, 59--63, 66--69, 104.]}
The following exact assertions are not proved or refuted here.
They are recorded separately from the proved uniqueness and
centralizer assertions.
\begin{enumerate}
\item For $U=\mathbb P^1_{\mathbb Q}\setminus\{0,1,\infty\}$,
every continuous section $s:G_{\mathbb Q}\to\pi_1(U)$ with
$C_\Pi(s(G_{\mathbb Q}^\circ))=1$ is the section of a rational
point. This is the smallest retained instance of the sufficiency
in scan 48 and of Conjecture A in scan 56.
\item Every continuous section $G_{\mathbb Q}\to G_{\mathbb Q(t)}$
has image in the decomposition group of a rational place.
By the checked return B7, this is exactly the assertion that its
geometric cyclotomic centralizer is nonzero. This is the smallest
retained birational instance of scan 45.
\item Every section of a hyperbolic curve admits a compatible
lift through the maps $\pi_1(U\setminus D)\to\pi_1(U)$ as the
finite closed subsets $D$ grow, sufficient to obtain a birational
section. The precise lifting problem is about continuous sections
and their geometric-conjugacy classes, not just individual lifts
of elements of $G_k$.
\item The assertion of Conjecture B' of scan 63 for arbitrary
abstract semilinear $G_{\mathbb Q}$-actions on the universal
pro-cover of $\mathbb P^1_{\overline{\mathbb Q}}
\setminus\{0,1,\infty\}$: every such action has a unique
$G_{\mathbb Q}^\circ$-fixed compatible geometric point in its
completed pro-cover. For actions with an open stabilizer for
every rational function, ZCF12 constructs an invariant finite
quotient over $\mathbb Q$ and reduces the assertion to
Conjecture B with its $G_k^\circ$ clause of scan 61.
Its fixed-point existence still contains (1),
but no additional finite-level obstruction is left in that
continuous case. The extra scope allowing abstract actions
without this continuity hypothesis remains unresolved.
\end{enumerate}
\end{result}
\begin{proof}[Verified deductions, attempts, and exact remaining steps]
For (1), (ZCF.2) proves that zero centralizer is equivalent to
noncuspidality. It proves neither existence of a point nor
existence of a fixed pro-point. The attached return D0b proves
at most one fixed pro-point, by passing to a characteristic finite
cover whose compactification has positive genus and applying
Jacobian Kummer injectivity to two putative fixed points. The
maps to that cover and then to its Jacobian are actual maps, so
this establishes uniqueness only. The existence of an interior
fixed pro-point is precisely the missing step.

We tried three further routes. First, Mochizuki's Hom theorem
ZCF1 requires open image; a point section has intersection $1$
with the infinite geometric kernel, and so its image is not open.
Second, the full topos theorem F9.2 supplies the exact map
\[
 \pi_0\operatorname{Hom}^{\mathrm{pin}}_{B G_k}
       (B G_k,U_{\mathrm{et}})
       \longrightarrow \operatorname{Sect}(E_U/G_k)/\Pi.
\]
Its source is $U(k)$ and its image is the point-section subset.
No construction in the proof lifts an arbitrary noncuspidal
section to the source. Establishing that lift for all such sections
would finish (1); the full-faithfulness theorem alone does not
produce it. Third, the finite-cover Kummer constraints in the
attached return D4 separate its specified two points via an explicit
degree-81 Fermat cover. They give compatibility requirements on
an arbitrary section at every finite level. They do not construct
a rational point from those requirements; that construction remains
unproved even for the displayed $U/\mathbb Q$.

For (2), B7 proves the exact equivalence, section by section,
between nonzero centralizer and a unique rational place. Its proof
uses a nonzero element $z$ and an open $U_0$ on which $z$ remains
nonzero. It then gets compatible unique cusps on every smaller
open. For an arbitrary section no such $z$ has been produced;
there is no first finite level to which that proof applies.
The $t$-adic sections show that the positive conclusion is realized,
but do not exhaust the sections. The plain full-$G_k$ centralizer
condition was refuted in ZCF2 and cannot supply the missing element.

For (3), the canonical surjection
$G_{k(U)}\twoheadrightarrow\pi_1(U)$ sends a decomposition-group
section at $x\in U(k)$ to the point section, and one at
$x\in\bar U\setminus U$ to a cuspidal section. This follows by
restricting covers to the henselian DVR and applying C2--C3.
The kernel of $\pi_1(U\setminus D)\to\pi_1(U)$ is the closed
normal subgroup generated by the inertia groups at $D$, by
normalization and purity for smooth curves. These facts specify
the quotients but do not give sections of their pullbacks by an
arbitrary $s$. Solving the finite embedding problems separately
would also leave a compatibility problem. This is the missing
cuspidalization step in the proposed ``forage de trous'' of scan 67.

For (4), ZCF12 supplies the previously missing comparison for
the explicit topology of pointwise convergence on the discrete
rational-function field. A finite orbit of a set of field
generators gives a $G_k$-stable finite quotient, and effective
descent gives a curve over $k$ itself. Thus, with this continuity
clause, Conjecture B with its $G_k^\circ$ strengthening of scan 61
and Conjecture B', quantified over all hyperbolic curves at the
same base field, are equivalent. Their existence conclusion has
not thereby been proved.

The remaining abstract-action problem has literal equations.
Fix the usual tangential continuous action $\rho_0$ on this
universal pro-cover and put
$A=\operatorname{Aut}_{\overline{\mathbb Q}}(\mathcal U)$.
Every abstract semilinear action is uniquely
$\rho(g)=a_g\rho_0(g)$ with
\[
 a_{gh}=a_g\rho_0(g)a_h\rho_0(g)^{-1}.
\]
We have neither constructed nor excluded such a cocycle whose
action normalizes none of the profinite deck groups of finite
hyperbolic quotients. The finite-orbit construction in ZCF12
excludes this possibility for every action with open function
stabilizers. Without that hypothesis the orbit of a field
generator can be infinite, and its generated field need not be
a finite covering level; this is the exact failed step of that
construction. No counterexample or substitute construction is
supplied here for those abstract actions. The further fixed-point
assertion in (4) remains unproved. ZCF7 refutes the additional
proposed extension to a cyclotomic ground field even for the
continuous actions coming from sections.

Thus the general reconstruction programme of scans 51 and 66--69
has these explicit remaining existence steps. F9.2 reconstructs
schemes from their full pinned \emph{\'etale topoi}; the displayed
forgetful map proves how this theorem relates to the smaller
fundamental-group input and locates exactly what is missing.
The text's predicate f), ``pas trop loufoque'', has no specified
mathematical definition and therefore supplies no further test.
\end{proof}

\begin{result}{ZCF7}{Refuted}{The proposed extension to cyclotomic ground fields}
\textup{[\S3, manuscript pp.\ 30--31; scans 62--63.]}
Let $F=\mathbb Q(\mu_\infty)$. The smooth projective normalization
$C/F$ of $z^6+2y^3+4=0$ has genus four, no $F$-rational point,
and a continuous section of $\pi_1(C)\twoheadrightarrow G_F$.
This refutes the claim that all the preceding existence statements
extend to fields finitely generated over a cyclotomic field.
\end{result}
\begin{proof}
This is the checked attached return D5 \cite{PriorCyclotomic}; we give the exact mechanism.
The map $(z,y)\mapsto(z^2:y:1)$ extends to a morphism to the
smooth cubic $D:X^3+2Y^3+4Z^3=0$. In a finite cyclotomic field
$\mathbb Q(\mu_N)$ the ramification degree $e$ at $2$ is a power
of two. With integral-normalized valuation, the three summands in
the cubic have residues $0,e,2e$ modulo three. The least nonzero
valuation is therefore attained once, which is incompatible with
their sum being zero; the same argument applies when a coordinate
vanishes. Hence $D(F)$ and then $C(F)$ are empty.
The degree-three map $C\to\mathbb P^1_z$ is ramified to index
three at the six simple zeros of $z^6+4$, and unramified at infinity.
Riemann--Hurwitz gives $2g-2=-6+12=6$.

D5 proves projectivity of $G_F$ by showing that every Brauer class
over every finite extension of $F$ is killed by a cyclotomic
extension, followed by the finite embedding-problem argument.
We checked that proof, including its reduction from finite modules
to Sylow subgroups and the Frattini-kernel induction. To apply it
here, take a cofinal descending sequence of characteristic open
subgroups $N_j$ of the finitely generated geometric group $\Pi_C$.
Starting at $E_C/\Pi_C=G_F$, projectivity lifts the section through
each finite-kernel quotient $E_C/N_{j+1}\to E_C/N_j$. The
compatible inverse limit is a continuous section of $E_C\to G_F$.
As $C$ is proper it has no cuspidal packet, and as $C(F)$ is empty
the section is not a point section. This is the actual extension
and the actual section-existence argument, not a deduction from a
section conjecture.
\end{proof}

\begin{result}{ZCF8}{Proved}{The commensurability assertion for universal algebraic pro-covers}
\textup{[\S3, manuscript p.\ 30; scan 62.]}
For smooth connected curves $U,V$ over an algebraically closed
field of characteristic zero, their universal connected finite-\'etale
pro-covers are isomorphic as schemes over that field exactly when
some connected finite \'etale cover of $U$ is isomorphic to a
connected finite \'etale cover of $V$. No compatibility of an
initial isomorphism with either deck group is assumed.
\end{result}
\begin{proof}
Write the universal pro-covers as inverse limits of their connected
finite covers; transition maps are finite, hence affine, and every
level is quasi-compact and quasi-separated. An isomorphism
$\widetilde U\simeq\widetilde V$ followed by projection to $U$
descends to a map $f:V_i\to U$, by the finite-presentation limit
theorem, Stacks Proposition 32.6.1 \cite{CF_StacksLimits}.
Using the inverse isomorphism, the induced map
$\widetilde U\to V_i$ descends to $g:U_j\to V_i$.
After increasing $j$, the equality on the inverse limit descends
to $fg=p_j:U_j\to U$, the finite \'etale transition map.
The map $f$ is nonconstant and surjective, since $fg$ is;
$g$ is nonconstant as well. For integral curves these maps are
quasi-finite. Moreover $g$ is proper: its graph is closed because
$V_i\to U$ is separated, and $U_j$ is proper over $U$.
Properness and quasi-finiteness also make $g$ finite. The image of $g$ is closed and dense in $V_i$, hence all of $V_i$.
Properness descends along this proper surjection: after any base
change, the image of a closed subset of $V_i$ is the image of its
closed inverse image in $U_j$. Thus $f$ is universally closed;
being separated and of finite type, it is proper. Quasi-finiteness
then makes $f$ finite.

At a geometric closed point the ramification indices for the
finite maps multiply. Their product for $fg$ is one, so the index
for $f$ is one at every point, since $g$ is surjective. Residue
extensions are separable in characteristic zero. A finite map of
smooth curves is flat and unramified when all these indices are
one, hence finite \'etale. Thus $V_i$ is the desired common cover.
Conversely, the covers of any connected finite \'etale cover $W$
are cofinal among covers of its base: intersect a subgroup of
$\pi_1(W)$ with its finitely many conjugates in the base group.
The resulting connected cover is finite \'etale over both. The
universal pro-cover of $W$ therefore identifies with the universal
pro-covers of both $U$ and $V$. This proves the converse and the
claimed equivalence of actual pro-schemes.
\end{proof}

\begin{result}{ZCF9}{Proved}{A basis of the more special elementary neighborhoods}
\textup{[\S3, manuscript pp.\ 33--34; scans 68--69.]}
Let $X$ be smooth over a characteristic-zero field $k$. Every point
has a fundamental system of Zariski-open neighborhoods which are
towers of hyperbolic curve fibrations, each curve fibration being
a finite \'etale cover of a relative punctured projective line
with finite \'etale boundary of degree at least three. The bottom
field may be a finite extension of $k$, accounting for geometric
connected components. This proves the geometric neighborhood
assertion of scan 68 in its characteristic-zero context; it does
not construct the additional sections needed in scan 67.
\end{result}
\begin{proof}
We work inside any prescribed affine open of $X$. The exact
Artin-neighborhood input is Panin--Stavrova--Vavilov, Proposition 1
\cite{CF_PSV2009}: over an infinite field, a smooth geometrically
irreducible variety and a finite set of closed points admit a
Zariski-open neighborhood with an elementary fibration to an open
of projective space of dimension one less. Definition 1 there
gives a smooth proper geometrically connected curve compactification
with nonempty finite \'etale boundary. Passing to the finite
constant field of a connected component gives the indicated bottom
field; in characteristic zero this component is geometrically
irreducible. For a nonclosed point choose a closed specialization
in the prescribed affine open; every open containing the
specialization also contains the original point.

It remains to prove the more special form of a single fibration,
retaining the chosen closed point $x$. Write
$C\to B$ for its smooth proper compactification and $D$ for its
nonempty finite \'etale boundary; let $b$ be the image of $x$.
On the smooth proper curve $C_b/k(b)$ choose $L=\mathcal O(nD_b)$
with $n$ large. There is a finite map $h_b:C_b\to\mathbb P^1$
defined by a pair of sections of $L$, with these properties:
its branch points are simple and have pairwise distinct images;
it is unramified at $x$ and at $D_b$; the images of $D_b$ are
pairwise distinct and avoid all branch values and $h_b(x)$;
and the entire fibre over $h_b(x)=0$ is unramified.

Here are details of this choice, to justify the retained-point
condition. Impose that the numerator vanish simply at the finite
reduced subscheme $x$ and that the denominator not vanish there.
For $n$ sufficiently large, Riemann--Roch makes sections with these
conditions surject onto arbitrary jets of order at most two at any
other point and first jets at any pair of distinct other points.
The pair has no common zero on a nonempty open of its parameter
space: a common zero imposes two independent conditions at a
variable point of a curve. Ramification of order at least three
imposes two conditions on its two-jet, also at a one-dimensional
variable point. Two distinct simple critical points with the same
value impose three independent conditions on their two first jets,
where the pair of points has dimension two. A critical point above
zero, or above the image of a point of $D_b$, similarly imposes
two conditions with only one variable point. At the finitely many
prescribed points the simple-zero, unramified, and distinct-value
conditions are nonempty open conditions on their jets and values.
These dimension counts use the specified surjections on jets;
their bad incidences have images of strictly smaller dimension.
They leave a nonempty open in the affine parameter space over
the infinite field $k(b)$, so a $k(b)$-defined pair exists.

For the same large $n$, proper cohomology and base change identifies
sections of $\mathcal O_C(nD)$ on the fibre with the fibre of a
locally free pushforward. Lift the chosen pair to a neighborhood
of $b$. Properness allows us to shrink that neighborhood so that
the pair has no common zero and defines a finite map
$h:C\to\mathbb P^1_B$. Its ramification divisor is finite \'etale
over $B$: its points were simple on the fibre at $b$, and the
failure locus is closed after taking its proper image in $B$.
The distinct branch images and the images of $D$ consequently
form disjoint finite \'etale divisors after further shrinking.
Add enough disjoint sections of $\mathbb P^1_B$ to make their
total degree at least three, avoiding $h_b(x)$ and all those
values at $b$, and shrink again. Their existence follows by
choosing finitely many distinct values in the infinite field
$k(b)$ and lifting their coordinates locally on $B$.

Let $V$ be the complement of this divisor in $\mathbb P^1_B$
and let $U'=h^{-1}(V)$. The map $U'\to V$ is finite \'etale;
$U'$ contains $x$ and lies in $C\setminus D$. The reduced
complement of $U'$ in $C$ is finite \'etale over $B$ as well.
At a simple ramification point the local equation is a unit
times $z^2$ after \'etale localization; its reduced inverse image
is the ramification section, disjoint from the unramified inverse
images. At the other removed values the inverse image is already
\'etale. This proves the last assertion without confusing a
reduced inverse image with the generally nonreduced branch fibre.
The fibre Euler characteristic is multiplied by the degree of
the finite \'etale cover, so it is negative.

Apply induction to a neighborhood of $b$ in the smooth base $B$
and pull $U'$ back to it. All changes are restrictions to actual
Zariski opens retaining the prescribed point; hence the final
tower is inside the initially prescribed open of $X$. This proves
the claimed fundamental-system assertion.
\end{proof}

\begin{result}{ZCF10}{Proved}{Generic cusps, closed-point determination, and higher-codimension decomposition groups}
\textup{[\S3, manuscript pp.\ 32, 36, 45--49; scans 65, 73,
91--99.]}
The following additional assertions used without proof in the
section argument hold under its hypotheses of a smooth integral
finite-type base in characteristic zero and a smooth relative
hyperbolic curve.
\begin{enumerate}
\item A rational point of the generic finite \'etale boundary
extends uniquely to a boundary section over the normal base.
\item Two rational sections which agree at all closed points
of a common nonempty open domain are equal as rational sections.
\item For $A=\mathcal O_{S,s}$, compatible henselizations inside
$K^{\mathrm{sep}}$, and $K^h=\operatorname{Frac}(A^h)$,
$K^{\mathrm{sh}}=\operatorname{Frac}(A^{\mathrm{sh}})$, the
decomposition and inertia groups are
\[
 N_s=G_{K^h},\qquad I_s=G_{K^{\mathrm{sh}}},\qquad
 1\to I_s\to N_s\to G_{k(s)}\to1.                         \tag{ZCF.3}
\]
The groups depend on the chosen prolongation by conjugation.
For an actual rational section $f_K$, its map $G_K\to E_U$
kills $I_s$ exactly when $f_K$ extends on a neighborhood of $s$.
The induced quotient section is then the section of $f(s)$.
\end{enumerate}
\end{result}
\begin{proof}
For (1), the connected component of the finite \'etale boundary
containing the generic point in question is finite birational over
the normal base, hence isomorphic to it: its finite coordinate
algebra lies in the fraction field and is integral over the
integrally closed base ring. This supplies the section.
For (2), the equalizer of two sections is closed because the
target is separated. Closed points are dense in every closed
subset of a finite-type scheme over a field. Thus its underlying
set is the entire common domain. The ideal defining the equalizer
vanishes in every prime and therefore is zero on the reduced
domain. The sections agree as morphisms, hence rationally.

For (3), normalize $S$ in finite Galois subextensions of
$K^{\mathrm{sep}}$. The stabilizer of a chosen point above $s$
acts on its residue field, with kernel inertia. Passing to the
limit, the unramified algebraic local extensions are precisely
the \'etale neighborhoods defining the henselization, and their
strictly local extensions are precisely those defining the strict
henselization. Hensel's lemma lifts each finite separable residue
extension, so the residue map is surjective. This proves every
term and map in (ZCF.3). The doubled tilde on the strict field
in scan 93 is mathematically necessary: using $K^h$ for both
groups would erase the nontrivial quotient $G_{k(s)}$.

If $f$ extends near $s$, functoriality and the strict local base
make its map on $I_s$ trivial. Conversely, base change to
$\operatorname{Spec}A^{\mathrm{sh}}$. The rational map to the
pulled-back curve has trivial map on the absolute group of its
fraction field: its image lies in the geometric fibre group, whose
map to $E_U$ is injective by C4. At each height-one point its inertia map is
therefore trivial. C2 extends the point over every such DVR;
finite presentation spreads these extensions and C1 extends
across the codimension-at-least-two complement. The extension
over $A^{\mathrm{sh}}$ is unique, agrees over the fraction field
with the given $K$-point, and descends by faithfully flat descent
to $A$. Finite presentation spreads it to a neighborhood of $s$.
C4 identifies the two arithmetic extensions for the henselian
base and its closed fibre: they have the same geometric kernel
and quotient $G_{k(s)}$. The descended section restricts to the
actual point $f(s)$ under this identification.

In scan 73 the map $G_K\to\pi_1(\mathcal O_{S,s})\to\pi_1(S)$
is surjective at both stages: a connected finite \'etale cover
of a normal integral scheme is normal integral, and its generic
and local restrictions remain integral and connected. This is
the finite-cover criterion for surjectivity. In scans 98--99,
composition with a decomposition-group section gives the finite
point section when $f$ extends, by the preceding argument; when
it meets the boundary it gives the tangential section calculated
in C3 and C9. The Tate factor in scan 91 is $T=\widehat{\mathbb Z}(1)$
with cyclotomic action. The word ``trivial'' cannot be interpreted
as a trivial $G_{k(s)}$-module over a number field: the actual
conjugation formula is $g\tau g^{-1}=\tau^{\chi(g)}$.
\end{proof}

\begin{result}{ZCF11}{Proved}{The exact meaning of the inner-action reformulation in scan 102}
\textup{[\S3, manuscript pp.\ 50--51; scans 101--104.]}
Let $1\to\Pi\to E\to G_K\to1$ be the arithmetic extension
of a hyperbolic curve and let $s:G_K\to E$ be a continuous
section. For an inertia group $T\subset G_K$, write
$\alpha_s:T\to\operatorname{Aut}(\Pi)$ for conjugation by $s$.
Triviality of its outer action is equivalent to the existence
of a continuous homomorphism $u_s:T\to\Pi$ with
$\alpha_s(\tau)=\operatorname{int}(u_s(\tau))$.
This homomorphism is unique. The outer condition is independent
of $s$, so one such section suffices and every such section then
has this property.
\end{result}
\begin{proof}
The map $\Pi\to\operatorname{Inn}(\Pi)$ is an isomorphism of
topological groups because $\Pi$ is centre-free (A3): it is a
continuous bijection from a compact group to a Hausdorff group.
Vanishing in $\operatorname{Out}(\Pi)$ places $\alpha_s(T)$
in this image and permits its unique continuous inverse. The
inverse preserves multiplication, which proves that $u_s$ is a
homomorphism. The converse follows by projecting inner
automorphisms to the outer group. Two sections differ by a
$\Pi$-valued crossed homomorphism, so their conjugation actions
differ by inner automorphisms. This proves the last assertion.
The actual map $s|_T:T\to E$ does not itself land in $\Pi$:
its composite with $E\to G_K$ is the inclusion of $T$.
The map $u_s$ just constructed is the precise comparison required
by the manuscript's phrase ``induit un hom $T\to\pi$''.

The suggestion that a criterion using only the Jacobian's
$p$-divisible group should detect good reduction of a proper
curve has the same obstruction as the abelian criterion:
the compact-type degeneration F17.2 has a Jacobian extending
to an abelian scheme, whose $p$-power kernels form an extending
Barsotti--Tate group for every $p$, while the curve is singular
in its stable reduction. The same counterexample exists with the
prime equal to the residue characteristic: take
$y^2=(x^3-5^6)(x^3-1)$ over $\mathbb Q_5$. The equations of F17.2
with $t=5$ give a stable model over $\mathbb Z_5$, with node
$VW=5$ and the two displayed elliptic components over $\mathbb F_5$.
Their discriminants are units at $5$, so both components are
smooth. The compact-type Jacobian extends to an abelian scheme
by the relative-Picard construction and descent checked in F17.2
\cite[\S1]{CF_DupuyRabinoff2017};
its kernels of multiplication by $5^n$ form an extending
Barsotti--Tate group, whereas the stable curve has a node.
The nonabelian criterion is F17.1.
The further construction of a nongeometric genuine action whose
outer action already is geometric, in scan 104, is the unresolved
section-existence problem ZCF6(1); F17.4 instead constructs an
outer action which is itself nongeometric.
\end{proof}

\begin{result}{ZCF12}{Proved}{Continuous semilinear actions descend to a finite quotient over the base field}
\textup{[\S3, manuscript pp.\ 29--33; scans 58--59, 62--63, 66--67.]}
Let $k$ have characteristic zero. A continuous section of
$E_U\to G_k$ for a hyperbolic curve $U/k$ defines a semilinear
action on its universal geometric pro-cover, equivariant for
the ordinary scalar action on $U_{\bar k}$. These are exactly
the continuous semilinear actions equivariant for that specified
projection.

More generally, let $\mathcal U$ be the universal pro-cover of
a hyperbolic curve over $\bar k$, and let $\Omega$ be the union
of the function fields of its finite covering levels. Give the
semilinear automorphism group the topology of pointwise convergence
on the discrete field $\Omega$. Every semilinear $G_k$-action
for which each $f\in\Omega$ has open stabilizer admits a finite
hyperbolic quotient $W/k$: there is an equivariant projection
$\mathcal U\to W_{\bar k}$, its deck group is normalized by
all of $G_k$, and the action comes from a continuous section of
$E_W\to G_k$. The geometric quotient can be chosen as a finite
\'etale cover of any prescribed initial level.

Consequently the family of Conjecture B with the $G_k^\circ$
strengthening of scan 61, and that of Conjecture B' at the same
field $k$, each quantified over all hyperbolic curves and with
this continuity clause in B', are equivalent. This statement
uses $G_k^\circ$-fixed compatible geometric pro-points on both
sides. Every semilinear
automorphism also extends canonically to the completed pro-cover.
The natural $G_k$-action on germs of geometric-conjugacy classes
of sections is well defined; the point-section map to those germs
is equivariant and injective for $k$ finitely generated over
$\mathbb Q$.
\end{result}
\begin{proof}
Fix a geometric fibre functor. The function field $\Omega$ of
the geometric universal cover of $U$ is also the union of the
finite arithmetic covering fields of $U/k$: a geometric cover
and its map descend to a finite coefficient extension, and every
arithmetic cover becomes a geometric cover after extending the
constants. Thus
$E_U=\operatorname{Aut}_{k(U)}(\Omega)$ with its Krull topology.
A section acts semilinearly on the directed system of finite
covers and hence on its inverse limit; a cofinal system of
finite arithmetic Galois covering fields is invariant.
The Krull topology and continuity of the section give open
stabilizers for every element of $\Omega$. Conversely an
action equivariant for the scalar action on $U_{\bar k}$ fixes
$k(U)$, giving a homomorphism $G_k\to E_U$ whose projection
is the identity. Open function stabilizers are exactly the
continuity condition for this homomorphism. Equivalently the
construction is the action on finite continuous $\Pi$-sets
and their fibre functor. This identifies actual sections after
the fibre functor has been fixed; changing it gives precisely
geometric inner conjugation.

We next prove the descent for an arbitrary action $\rho$ with
the stated topology. Start with any level $U_0/\bar k$ and
choose field generators
$L_0=\bar k(U_0)=\bar k(f_1,\ldots,f_m)\subset\Omega$.
An open subgroup of the profinite group $G_k$ has finite index.
Hence each orbit $\{\rho(g)f_j:g\in G_k\}$ is finite. Put
\[
 L=\bar k\bigl(\rho(g)f_j:
                   g\in G_k,\ 1\le j\le m\bigr),\qquad
 P=\operatorname{Aut}_L(\Omega).                         \tag{ZCF.4}
\]
The extension $L/L_0$ is finite: its finitely many generators
lie in the algebraic extension $\Omega/L_0$. It corresponds
to a connected finite \'etale cover $W_{\bar k}\to U_0$,
because $\Omega$ is the union of precisely those covering
fields. More explicitly, put all the generators in one finite
Galois \'etale covering field; its intermediate field $L$ is
the function field of the finite \'etale quotient by the
corresponding subgroup. Thus $P$ is an open subgroup of the
initial deck group, and $\mathcal U\to W_{\bar k}$ is its
universal geometric pro-cover. The field $L$ is $G_k$-stable,
so $\rho(g)P\rho(g)^{-1}=P$ for every $g\in G_k$.

This normalization descends the action on the actual open
curve, not merely its function field. Write
$p:\mathcal U\to W_{\bar k}$. Regarding a semilinear map as
a linear map to the appropriate scalar-conjugate curve,
$p\rho(g)$ factors through a finite covering level by
Stacks, Proposition 32.6.1 \cite{CF_StacksLimits}.
Refine that level to a finite Galois cover $W'\to W_{\bar k}$
corresponding to an open normal subgroup $P'\subset P$.
The descended map $h:W'\to g^*W_{\bar k}$ is invariant
under $P/P'$: for $a\in P$ the equality
$p\rho(g)a=p(\rho(g)a\rho(g)^{-1})\rho(g)=p\rho(g)$
holds, and the projection from $\mathcal U$ to $W'$ is
faithfully flat, being an inverse limit of finite \'etale
surjections.
Finite faithfully flat descent along $W'\to W_{\bar k}$
therefore gives a morphism $W_{\bar k}\to g^*W_{\bar k}$.
In affine descent charts the descended coordinate functions
are exactly the invariants in the finite Galois algebra;
uniqueness glues these morphisms. The inverse action gives
its inverse, and the action law follows after pullback along
$p$. We have therefore obtained the full semilinear action
on the open curve. It extends uniquely to its smooth proper
compactification and preserves its boundary.

For completeness we make the arithmetic descent effective.
The finite-generator argument in the checked prior return D11
\cite{PriorCyclotomic} works for any prescribed level as follows.
Choose a finite coefficient extension $k_0/k$ over which the
smooth compactification, its boundary, field generators, and
their finite defining relations are defined. Intersect
$G_{k_0}$ with the stabilizers of those generators and take
the normal core in $G_k$. The result is an open normal subgroup
$H=G_{k_i}$. It fixes the generators and acts on the coefficients
by their usual scalar action, so the level has a $k_i$-form
on which this is exactly the standard $H$-action. Equality on
the function field gives equality of the projection maps:
one can factor them through a common finite level and use
integrality of that level and separatedness of the target.
In particular the projection is equivariant, including its
boundary. This proves directly the prior return's claim for
every initial level after an open restriction.

Apply this construction to the $G_k$-stable level just obtained.
Its field of $H$-invariants is $k_i(W_i)$. Normality of $H$
implies that every $\rho(g)$ preserves that field, with its
scalar action on $k_i$. Thus the finite group
$G_k/H=\operatorname{Gal}(k_i/k)$ acts semilinearly on the
smooth proper marked pair $(\bar W_i,D_i)$ over $k_i$.
The action is an actual Galois descent datum. It is effective:
$\omega_{\bar W_i}(D_i)^{\otimes3}$ is very ample for a
hyperbolic marked curve, is canonically linearized by every
isomorphism of the pair, and embeds both the curve and its
marked divisor into the projective space of its sections.
Finite Galois vector-space descent identifies that section
space with $k_i\otimes_k$ its invariant space; the trace-dual
basis proof and the descent of all homogeneous ideal pieces
are exactly those written in F17.7. Descending these ideals
gives a projective pair $(\bar W,D)/k$ whose base change is
the displayed pair. Smoothness and the finite \'etale property
of $D$ descend by faithful flatness. Its complement $W$ is
hyperbolic because its geometric curve is a finite \'etale
cover of the initial hyperbolic curve. This proves the claimed
quotient over $k$ itself.

The field $\Omega$ is the union of the finite arithmetic
covering fields of this $W$, by the same coefficient descent
used in the first paragraph. The action fixes $k(W)$ and has
its prescribed action on $\bar k$. It therefore gives the
actual continuous section in
\[
 1\longrightarrow P\longrightarrow
 \operatorname{Aut}_{k(W)}(\Omega)=E_W
 \longrightarrow G_k\longrightarrow1.                    \tag{ZCF.5}
\]
Continuity again follows from the open stabilizers of finite
sets of field elements. This proves that every action in the
stated continuous form of B' is an ordinary section action
for a quotient curve over the same field. Conversely every
section action satisfies that continuity clause.

The comparison of completed geometric pro-points is explicit.
The pointed finite covers of $W_{\bar k}$ form a cofinal
subsystem of the pointed finite covers of $U_0$: composition
gives a cover of $U_0$, and the component through the chosen
lift in $V\times_{U_0}W_{\bar k}$ refines any cover
$V\to U_0$. Extending the maps to proper compactifications
therefore gives a bijection
\[
 \varprojlim_{V\to U_0}\bar V(\bar k)
 \ \xleftarrow{\ \sim\ }\quad
 \varprojlim_{V\to W_{\bar k}}\bar V(\bar k).
\]
In group notation, if $\Delta_0$ is the original deck group,
the opens $M\subset P$ are cofinal among opens $N\subset
\Delta_0$, by taking $M\subset N\cap P$. The inverse
bijection sends the coordinate at such an $M$ to the level
for $N$; common refinement proves independence of this choice.
It is $G_k$-equivariant: both actions extend the same action
on $\mathcal U$, and extension of every finite-level map to
its proper compactification is unique. Hence the two assertions
about a unique $G_k^\circ$-fixed compatible geometric pro-point
concern exactly the same actions and points. This proves their
fixed-base equivalence without proving either existence assertion.
For a general characteristic-zero field no replacement of
$G_k^\circ$ by $G_k$ is part of this argument. Over a field
finitely generated over $\mathbb Q$, the checked uniqueness
statement D0b \cite{PriorCyclotomic} supplies the additional
comparison: a unique $G_k^\circ$-fixed point is $G_k$-fixed by
normality, while a $G_k$-fixed point is the unique
$G_k^\circ$-fixed point. In that arithmetic range the plain
existence formulation of B has the same fixed-point problem.

Every finite map between smooth curves extends uniquely between
their smooth proper compactifications, by the DVR valuative
criterion for the proper target. Hence transition maps and the
action extend compatibly to the inverse limit of compactifications.
For an isomorphism of universal pro-covers without a specified
deck group, the descent argument of ZCF8 supplies finite-level
comparison maps in both directions; they extend to the proper
compactifications, and their inverse-limit composites are the
identity. This proves the claimed canonical extension of such an
isomorphism. For a semilinear isomorphism, apply the same argument
to the corresponding scalar-conjugate curve.

For the assertion on germs, a section over an open $H\subset G_k$
is identified with its restrictions to smaller open subgroups and
with its conjugates by $\Pi$. For $g\in G_k$, choose $e\in E_U$
above $g$ and send $s$ to
$h\mapsto e\,s(g^{-1}hg)\,e^{-1}$ on $gHg^{-1}$.
Another choice of $e$ differs by $\Pi$ and gives the same class;
restriction commutes with this operation. The group law is valid
modulo exactly the same conjugation, so it is an action on germs.
A geometric point $x\in U(\bar k)$ is defined over a finite
extension; its section defines a germ, and functoriality gives
equivariance. Equality of two such germs is equality of the two
point-section classes over a common finite extension. B5 then
implies equality of the points. The assertion that the noncuspidal
germs are all in this image is precisely the unresolved surjectivity
in ZCF6, not a consequence of this equivariant injection.
\end{proof}

